\subsection{Factorización Rápida de Enteros Grandes (Pollard's Rho + Miller-Rabin)}

\graybold{Preguntas guía:}

\begin{itemize}
\item ¿Necesito factorizar en primos números de hasta $\sim 10^{12}$ o más, donde la división de prueba ($O(\sqrt{n})$) es demasiado lenta con muchos números o primos grandes?
\item ¿Necesito verificar primalidad de números grandes rápido y con certeza práctica?
\item ¿El problema involucra propiedades de los factores primos de varios números (compatibilidad, MCD, etc.)?
\end{itemize}

\graybold{¿Para qué sirve?} Factorizar números grandes (hasta $10^{12}$ o más) en tiempo esperado muy por debajo de $O(\sqrt{n})$ por número.

\graybold{Complejidad:} Miller-Rabin: \bigO{\log^3 n} por chequeo. Pollard's Rho: \bigO{n^{1/4}} esperado por factor encontrado.

\graybold{Fundamento teórico:} Miller-Rabin es un test de primalidad basado en las propiedades de las raíces cuadradas de 1 módulo un primo; con un conjunto fijo de bases testigo es determinístico en el rango de este problema. Pollard's Rho encuentra un factor no trivial de $n$ simulando una secuencia pseudoaleatoria $x_{i+1}=x_i^2+c \pmod n$ y buscando una colisión módulo un factor de $n$ (mucho más frecuente que módulo $n$ completo), detectada vía el algoritmo de ciclos de Floyd combinado con mcd. Factorizar completamente un número se reduce a aplicar Pollard's Rho recursivamente sobre los factores compuestos que va encontrando.

\graybold{Ejercicio aplicado}

(Primera mitad de "Holy Network", ver también la entrada de Clique Máximo en Algoritmos de Grafos) Dados N números de hasta $10^{12}$, obtener el conjunto de primos que divide a cada uno.

\graybold{I/O:} N ≤ 50 por caso → no hace falta ninguna optimización de lectura (patrón 1); el costo está en la factorización, no en la E/S. Verificado: división de prueba tarda \textasciitilde{}5.6s en un caso adversario donde Pollard's Rho tarda \textasciitilde{}0.09s.

\graybold{Código solución}

\begin{lstlisting}[language=Python]
import random
from math import gcd

def is_prime(n):
    if n < 2:
        return False
    for p in (2, 3, 5, 7, 11, 13, 17, 19, 23, 29, 31, 37):
        if n % p == 0:
            return n == p
    d = n - 1
    r = 0
    while d % 2 == 0:
        d //= 2
        r += 1
    for a in (2, 3, 5, 7, 11, 13, 17, 19, 23, 29, 31, 37):
        x = pow(a, d, n)
        if x == 1 or x == n - 1:
            continue
        for _ in range(r - 1):
            x = x * x % n
            if x == n - 1:
                break
        else:
            return False        # testigo compuesto encontrado
    return True

def pollard_rho(n):
    if n % 2 == 0:
        return 2
    while True:                 # reintenta con otro c si falla
        c = random.randint(1, n - 1)
        f = lambda x: (x * x + c) % n
        x = y = random.randint(2, n - 1)
        d = 1
        while d == 1:
            x = f(x)
            y = f(f(y))          # tortuga y liebre (deteccion de ciclos)
            d = gcd(abs(x - y), n)
        if d != n:               # d==n: colision trivial, reintentar
            return d

def prime_factors(n):
    if n == 1:
        return set()
    if is_prime(n):
        return {n}
    d = pollard_rho(n)
    return prime_factors(d) | prime_factors(n // d)
\end{lstlisting}


\subsection{Clique Máximo en Grafos Pequeños (Bron-Kerbosch con Poda)}

\graybold{Preguntas guía:}

\begin{itemize}
\item ¿Piden el subconjunto más grande de elementos tal que TODOS los pares cumplan una relación (compatibilidad, conflicto, etc.)?
\item ¿El número de elementos es pequeño (hasta $\sim$50-100), aunque el problema en general sea NP-difícil?
\item ¿La relación de compatibilidad no tiene una estructura especial (bipartita, de intervalos) que permita un atajo polinomial?
\end{itemize}

\graybold{¿Para qué sirve?} Encontrar el clique máximo en un grafo pequeño, cuando no hay estructura especial que permita resolverlo en tiempo polinomial (como sí ocurre en grafos bipartitos con el teorema de König).

\graybold{Complejidad:} Exponencial en el peor caso (NP-difícil), pero con poda por pivote resulta muy rápido en la práctica para $n \le \sim 50$-100 (verificado: <10ms incluso en grafos aleatorios densos $G(50, 0.5)$).

\graybold{Fundamento teórico:} Bron-Kerbosch explora el espacio de cliques manteniendo tres conjuntos: $R$ (clique parcial), $P$ (candidatos) y $X$ (descartados). En cada paso se elige un "pivote" (el nodo con más vecinos en $P$) y solo se ramifica sobre los nodos que NO son vecinos del pivote. Para clique MÁXIMO se agrega una poda: si el tamaño parcial más los candidatos restantes no supera el mejor encontrado, se corta la rama.

\graybold{Ejercicio aplicado}

Dados N ingredientes con un número asociado (hasta $10^{12}$), donde dos son compatibles si comparten un factor primo mayor a 1, hallar el tamaño del grupo más grande donde todos los pares son compatibles entre sí.

\graybold{I/O:} N ≤ 50 → entrada trivial; toda la complejidad está en la factorización (entrada anterior) y en la búsqueda de clique.

\graybold{Código solución}

\begin{lstlisting}[language=Python]
import sys
import random
from math import gcd

def is_prime(n):
    if n < 2:
        return False
    for p in (2, 3, 5, 7, 11, 13, 17, 19, 23, 29, 31, 37):
        if n % p == 0:
            return n == p
    d = n - 1
    r = 0
    while d % 2 == 0:
        d //= 2
        r += 1
    for a in (2, 3, 5, 7, 11, 13, 17, 19, 23, 29, 31, 37):
        x = pow(a, d, n)
        if x == 1 or x == n - 1:
            continue
        for _ in range(r - 1):
            x = x * x % n
            if x == n - 1:
                break
        else:
            return False
    return True

def pollard_rho(n):
    if n % 2 == 0:
        return 2
    while True:
        c = random.randint(1, n - 1)
        f = lambda x: (x * x + c) % n
        x = y = random.randint(2, n - 1)
        d = 1
        while d == 1:
            x = f(x)
            y = f(f(y))
            d = gcd(abs(x - y), n)
        if d != n:
            return d

def prime_factors(n):
    if n == 1:
        return set()
    if is_prime(n):
        return {n}
    d = pollard_rho(n)
    return prime_factors(d) | prime_factors(n // d)

def max_clique_size(n, adj):
    # adj[i] = bitmask (entero) de los vecinos de i
    best = 0
    full_mask = (1 << n) - 1

    def bronkerbosch(r_size, P, X):
        nonlocal best
        if P == 0 and X == 0:
            best = max(best, r_size)
            return
        if r_size + bin(P).count('1') <= best:
            return                       # poda: no puede mejorar el maximo
        PX = P | X
        pivot, max_count = -1, -1
        temp = PX
        while temp:
            v = (temp & -temp).bit_length() - 1
            temp &= temp - 1
            cnt = bin(adj[v] & P).count('1')
            if cnt > max_count:
                max_count, pivot = cnt, v
        candidates = P & ~adj[pivot]     # solo NO-vecinos del pivote
        while candidates:
            v = (candidates & -candidates).bit_length() - 1
            candidates &= candidates - 1
            bronkerbosch(r_size + 1, P & adj[v], X & adj[v])
            P &= ~(1 << v)
            X |= (1 << v)

    bronkerbosch(0, full_mask, 0)
    return best

def solve():
    data = sys.stdin.buffer.read().split()
    idx = 0
    out = []
    while idx < len(data):
        n = int(data[idx]); idx += 1
        if n == 0:
            break
        vals = [int(data[idx + i]) for i in range(n)]
        idx += n

        factor_sets = [prime_factors(v) for v in vals]

        adj = [0] * n
        for i in range(n):
            for j in range(i + 1, n):
                if factor_sets[i] & factor_sets[j]:   # comparten un primo
                    adj[i] |= (1 << j)
                    adj[j] |= (1 << i)

        out.append(str(max_clique_size(n, adj)))
    sys.stdout.write("\n".join(out) + "\n")

solve()
\end{lstlisting}